% history-id: range-example
\section{Zasięg i przykład}

Dla rzutu z poziomu gruntu czas lotu wynika z warunku \(y(T)=0\). Po
odrzuceniu rozwiązania \(T=0\) otrzymujemy
\[
T=\frac{2v_0\sin\alpha}{g}.
\]

Zasięg poziomy jest równy
\[
R=v_0\cos\alpha\,T
 =\frac{v_0^2\sin(2\alpha)}{g}.
\]

\subsection{Kąt maksymalnego zasięgu}

Dla ustalonej wartości \(v_0\) oraz stałego \(g\) o zasięgu decyduje czynnik
\(\sin(2\alpha)\). Jego największa wartość jest równa \(1\), zatem
\[
2\alpha=90^\circ
\quad\Longrightarrow\quad
\alpha=45^\circ.
\]

Ponadto kąty dopełniające się do \(90^\circ\) dają w tym modelu taki sam
zasięg, ponieważ
\[
\sin(2\alpha)=\sin\left(2(90^\circ-\alpha)\right).
\]
Nie oznacza to jednak identycznych torów: większy kąt daje wyższy lot i
dłuższy czas przebywania w powietrzu.

\subsection{Przykład liczbowy}

Niech
\[
v_0=20\,\mathrm{m/s},
\qquad
\alpha=45^\circ,
\qquad
g=9{,}81\,\mathrm{m/s^2}.
\]

Czas lotu wynosi
\[
T=\frac{2\cdot20\sin45^\circ}{9{,}81}
 \approx 2{,}88\,\mathrm{s}.
\]

Maksymalna wysokość:
\[
H=\frac{20^2\sin^2 45^\circ}{2\cdot9{,}81}
 \approx 10{,}2\,\mathrm{m}.
\]

Zasięg:
\[
R=\frac{20^2\sin90^\circ}{9{,}81}
 \approx 40{,}8\,\mathrm{m}.
\]

Te trzy liczby opisują najważniejsze cechy lotu w modelu bez oporu powietrza:
jak długo ciało pozostaje w ruchu, jak wysoko się wznosi i jak daleko dociera.
